% =============================================================================
% БИЛЕТ 7
% =============================================================================

\ticket{7}{}

\question{Дифференцируемое отображение, производная отображения. Необходимое условие дифференцируемости обратного отображения, достаточное условие дифференцируемости обратного отображения}
\begin{defn}[Дифференцируемое отображение]
  По определению, это $F(x) - F(x^0) = A(x-x^0) + \varepsilon(x)\|x-x^0\|$, где $A$ --- матрица Якоби, а $\varepsilon(x) \to 0$ при $x \to x^0$. Матрица Якоби --- это матрица, которая состоит из всех частных производных. И если $F$ дифференцируемо в точке $x^0$, то матрица Якоби --- это производная отображения $F$ в точке $x^0$.
\end{defn}

\begin{note}
Например для отображения $F : R^n \to R^m$ матрица Якоби --- это матрица размера $m \times n$, в которой на пересечении $i$-ой строки и $j$-ого столбца стоит $\dfrac{\partial F_i}{\partial x_j}$.
\end{note}
\begin{theorem}[О производной обратного отображения]
Пусть $\Omega \subset \R^n$ --- открытое множество и $x^0$ --- внутренняя точка $\Omega$. Пусть
\begin{enumerate}
  \item Отображение $F : \Omega \to \R^n$ обратимо на $\Omega$;
  \item Отображение $F$ дифференцируемо в точке $x^0$;
  \item Отображение $F^{-1} : F(\Omega) \to \R^n$ непрерывно в точке $y^0 = F(x^0)$;
  \item Якобиан не равен нулю: $\mathcal{J}_F(x^0) \neq 0$;
  \item Точка $y^0$ --- внутренняя точка $F(\Omega)$.
\end{enumerate}
Тогда $F^{-1}$ дифференцируемо в точке $y^0$ и
\[
  \mathcal{D}F^{-1}(y^0) = \bigl(\mathcal{D}F(x^0)\bigr)^{-1}.
\]
\begin{proof}
Пусть $A$ --- матрица Якоби $F$ в точке $x^0$. Тогда по определению
дифференцируемости $F$:
\[
  F(x) - F(x^0) = A(x - x^0) + \varepsilon(x)\|x - x^0\|,
  \quad \varepsilon(x) \to 0 \text{ при } x \to x^0.
\]
Так как $F$ обратимо, подставим $x = F^{-1}(y)$, $x^0 = F^{-1}(y^0)$,
тогда $F(x) - F(x^0) = y - y^0$:
\[
  y - y^0
  = A\bigl(F^{-1}(y) - F^{-1}(y^0)\bigr)
  + \varepsilon(F^{-1}(y))\,\|F^{-1}(y) - F^{-1}(y^0)\|.
\]
Так как $\mathcal{J}_F(x^0) \neq 0$, матрица $A$ обратима.
Умножим обе части на $A^{-1}$ слева:
\[
  F^{-1}(y) - F^{-1}(y^0)
  = A^{-1}(y - y^0)
  - A^{-1}\varepsilon(F^{-1}(y))\,\|F^{-1}(y) - F^{-1}(y^0)\|.
\]
Покажем, что второе слагаемое есть $o(\|y - y^0\|)$.
Перепишем его как $\tilde{\varepsilon}(y)\cdot\|y-y^0\|$, где
\[
  \tilde{\varepsilon}(y)
  = -A^{-1}\varepsilon(F^{-1}(y))\cdot
  \frac{\|F^{-1}(y) - F^{-1}(y^0)\|}{\|y - y^0\|}.
\]
Это произведение двух множителей.

\textit{Первый множитель:} $\varepsilon(F^{-1}(y)) \to 0$ при $y \to y^0$
по теореме о пределе композиции: из непрерывности $F^{-1}$ в $y^0$ следует
$F^{-1}(y) \to x^0$, а $\varepsilon(x) \to 0$ при $x \to x^0$.

\textit{Второй множитель:} по основной геометрической теореме,
применённой к $F$ в точке $x^0$:
\[
  \|y - y^0\| = \|F(x) - F(x^0)\| \geq c\,\|F^{-1}(y) - F^{-1}(y^0)\|,
\]
откуда $\dfrac{\|F^{-1}(y) - F^{-1}(y^0)\|}{\|y - y^0\|} \leq \dfrac{1}{c}$ --- ограничен.

Бесконечно малое на ограниченное даёт бесконечно малое,
поэтому $\tilde{\varepsilon}(y) \to 0$ при $y \to y^0$. Итого:
\[
  F^{-1}(y) - F^{-1}(y^0)
  = A^{-1}(y - y^0) + \tilde{\varepsilon}(y)\|y - y^0\|,
  \quad \tilde{\varepsilon}(y) \to 0 \text{ при } y \to y^0.
\]
Это в точности определение дифференцируемости $F^{-1}$ в точке $y^0$,
и его производная равна $\mathcal{D}F^{-1}(y^0) = A^{-1} = \bigl(\mathcal{D}F(x^0)\bigr)^{-1}$.
\end{proof}
\end{theorem}
\question{Степенные ряды с комплексными членами. Круг и радиус сходимости. Характер сходимости в круге. Формула Коши-Адамара}

% Ответ...

